% Paper 2 Abstract
\begin{abstract}
Uniform quantization of the KV cache---assigning the same bit-width to every dimension---is provably suboptimal when the variance across dimensions is non-uniform. We show this non-uniformity is extreme in practice: key cache vectors exhibit AM/GM ratios up to 57$\times$ on outlier layers, corresponding to a coding gain of +2.9 bits (4-bit adaptive $\approx$ 7-bit uniform quality). We present \textbf{PCA-Quant}, a framework that rotates key vectors into their PCA basis and applies water-filling bit allocation---giving more bits to high-variance dimensions and fewer to low-variance ones. We prove PCA maximizes coding gain over all orthogonal transforms (including the Hadamard rotation used by TurboQuant), and show that the optimal tiered allocation follows a simple monotonicity rule solvable in closed form.

Empirical characterization across three models reveals that (1) coding gain concentrates on outlier layers that cause catastrophic quantization failure, making PCA-Quant a targeted fix rather than a uniform improvement; (2) V cache is near-isotropic (AM/GM $\approx 1.4$), so adaptive allocation is needed only for K; and (3) PCA bases are stable across sequence lengths for normal layers (CV $<$ 10\%), enabling one-time offline calibration.

We further investigate combining PCA-Quant with attention merging (AM), a token-level compression method. A surprising finding: AM at $2\times$ compression \emph{improves} reading comprehension accuracy by +2.2\% over the uncompressed baseline---a ``compression as regularization'' effect---followed by a sharp cliff between $2\times$ and $5\times$. This cliff-plateau pattern challenges the assumption that compression monotonically degrades quality and motivates quality-curve evaluation rather than single-point reporting.
\end{abstract}
